% gauge-filter.tex — Gauge filter: raw objects → G_T barrier → invariant structure
% Requires opoch.sty (loaded by main document)

\begin{figure}[htbp]
\centering
\begin{tikzpicture}[
  obj/.style={rounded rectangle, draw=#1, fill=#1!10, minimum width=10mm,
    minimum height=7mm, font=\small, inner sep=2pt},
  obj/.default={opochgray},
  blocked/.style={rounded rectangle, draw=opochred!50, fill=opochred!8,
    minimum width=10mm, minimum height=7mm, font=\small, inner sep=2pt,
    text=opochred!60},
  passed/.style={rounded rectangle, draw=opochgreen, fill=opochgreen!10,
    minimum width=10mm, minimum height=7mm, font=\small\bfseries, inner sep=2pt,
    text=opochgreen!80!black},
  barrier/.style={draw=opochorange, fill=opochorange!8, line width=1.2pt,
    minimum width=8mm, minimum height=50mm, rounded corners=2pt},
  arr/.style={-{Stealth[length=2.5mm]}, opochgray!60, thick},
  xarr/.style={-{Stealth[length=2mm]}, opochred!50, thick, dashed},
  garr/.style={-{Stealth[length=2.5mm]}, opochgreen!70, thick},
]

  % === Left side: Raw encoding-dependent objects ===
  \node[font=\footnotesize\bfseries, opochgray] at (-4.2, 2.8) {Raw objects};

  % Passing objects (encoding-invariant)
  \node[obj=opochblue] (r1) at (-4.2,  1.8) {$3$};
  \node[obj=opochblue] (r2) at (-4.2,  0.6) {$\pi$};
  \node[obj=opochblue] (r3) at (-4.2, -0.6) {$e$};

  % Blocked objects (encoding-dependent labels)
  \node[blocked] (b1) at (-4.2, -1.8) {\textsf{``three''}};
  \node[blocked] (b2) at (-4.2, -3.0) {\textsf{``drei''}};
  \node[blocked] (b3) at (-4.2, -4.2) {\textsf{``III''}};

  % === Barrier: Gauge filter G_T ===
  \node[barrier] (filt) at (0, -0.7) {};
  \node[font=\small\bfseries, opochorange!80!black, rotate=90] at (0, -0.7)
    {$\gauge$ filter};

  % Dashed vertical lines for barrier texture
  \foreach \y in {-2.8,-2.2,-1.6,-1.0,-0.4,0.2,0.8,1.4} {
    \draw[opochorange!30, thin] (-0.25, \y) -- (0.25, \y);
  }

  % === Arrows: passing through ===
  \draw[arr] (r1) -- (-0.6,  1.8);
  \draw[arr] (r2) -- (-0.6,  0.6);
  \draw[arr] (r3) -- (-0.6, -0.6);

  \draw[garr] (0.6,  1.8) -- (3.0,  1.8);
  \draw[garr] (0.6,  0.6) -- (3.0,  0.6);
  \draw[garr] (0.6, -0.6) -- (3.0, -0.6);

  % === Arrows: blocked ===
  \draw[xarr] (b1) -- (-0.6, -1.8);
  \draw[xarr] (b2) -- (-0.6, -3.0);
  \draw[xarr] (b3) -- (-0.6, -4.2);

  % X marks for blocked items
  \node[font=\large\bfseries, opochred!60] at (-0.6, -1.8) {$\times$};
  \node[font=\large\bfseries, opochred!60] at (-0.6, -3.0) {$\times$};
  \node[font=\large\bfseries, opochred!60] at (-0.6, -4.2) {$\times$};

  % === Right side: Gauge-invariant structure ===
  \node[font=\footnotesize\bfseries, opochgreen!70!black] at (4.5, 2.8)
    {Gauge-invariant};

  \node[passed] (p1) at (4.5,  1.8) {$3$};
  \node[passed] (p2) at (4.5,  0.6) {$\pi$};
  \node[passed] (p3) at (4.5, -0.6) {$e$};

  \draw[garr] (3.0,  1.8) -- (p1);
  \draw[garr] (3.0,  0.6) -- (p2);
  \draw[garr] (3.0, -0.6) -- (p3);

  % Annotation
  \node[font=\scriptsize, opochgray, align=center, text width=55mm] at (0, -5.5)
    {The gauge filter $\gauge$ blocks encoding-dependent labels\\
     and passes only structure invariant under all\\
     admissible re-encodings.};

\end{tikzpicture}
\caption{The gauge filter~$\gauge$. Raw objects enter from the left.
  Encoding-invariant quantities ($3$, $\pi$, $e$) pass through; encoding-dependent
  labels (``three'', ``drei'', ``III'') are blocked. Only gauge-invariant
  structure survives as admissible truth.}
\label{fig:gauge-filter}
\end{figure}
